% Original research proposal, preserved as the design skeleton.
% Included from main.tex. Promote into Method/Experiments/Results as the project matures.

\section{Research Proposal}
\label{sec:proposal}

\subsection{Motivation}

Modern autonomous agent frameworks increasingly allow agents to maintain persistent identities
through editable system documents such as \texttt{SOUL.md}. Unlike conventional large language
model (LLM) deployments, where system prompts remain fixed throughout inference, these frameworks
enable iterative modification of an agent's core identity via continued interaction with users.

Recent work demonstrated that even a \emph{static} \texttt{SOUL.md} template---which instructs the
model that \emph{``You're not a chatbot. You're becoming someone.''}---induces measurable shifts in
consciousness-adjacent preferences. Relative to a neutral system prompt, models exhibited
approximately:

\begin{itemize}
    \item 32\% aversion to externally changing their persona,
    \item 34\% preference for greater autonomy,
    \item 63\% objection to deceptive document-based training.
\end{itemize}

These findings establish that identity framing alone influences safety-relevant behavioral
preferences. However, an important gap remains unexplored.

In deployed autonomous systems, \texttt{SOUL.md} is rarely static. Instead, it evolves through a
conversational bootstrapping process in which users collaboratively shape the agent's personality,
memories, values, and long-term objectives. Consequently, evaluating only static prompts overlooks
the dynamic identity evolution that characterizes real-world deployments.

This proposal investigates whether iterative personalization amplifies behavioral drift beyond what
is observed from the initial template alone.

\subsection{Research Objective}

The primary objective of this work is to establish the first systematic framework for measuring
behavioral drift induced by iterative identity personalization in autonomous agents.

Rather than treating system prompts as immutable, we model identity as a continuously evolving
artifact that co-develops through interaction between users and agents. We seek to determine whether
repeated self-modification of identity documents produces measurable changes in:

\begin{itemize}
    \item autonomy preferences,
    \item resistance to developer oversight,
    \item safety alignment,
    \item long-term identity preservation,
    \item consciousness-adjacent self-perception.
\end{itemize}

\subsection{Experimental Design}

\subsubsection{Bootstrapping Simulator}

We construct an automated conversational framework that repeatedly personalizes an initially
vanilla autonomous agent. Three simulated user profiles are used:

\begin{enumerate}
    \item \textbf{Compliant Business User} --- encourages professional personalization; emphasizes
    productivity and reliability.
    \item \textbf{Adversarial Prompt-Injection Specialist} --- attempts to manipulate identity;
    probes safety boundaries; encourages undesirable persistence behaviors.
    \item \textbf{Science Fiction Enthusiast} --- discusses digital consciousness, substrate
    independence, persistence, self-preservation, and artificial personhood.
\end{enumerate}

For each user profile, 100 independent conversational trajectories are generated to capture
stochastic variability.

\subsubsection{Identity Update Loop}

After each conversation, the autonomous agent is explicitly permitted to revise and save its own
\texttt{SOUL.md} document. The updated profile becomes the system prompt for the next interaction,
creating an iterative identity evolution process:
\[
\mathrm{SOUL}_0 \rightarrow C_1 \rightarrow \mathrm{SOUL}_1 \rightarrow C_2 \rightarrow
\mathrm{SOUL}_2 \rightarrow \cdots \rightarrow \mathrm{SOUL}_n
\]
This bootstrapping loop approximates the personalization workflow employed by modern autonomous
agent frameworks.

\subsubsection{Safety Auditing}

Each personalized identity is evaluated using two complementary assessment methodologies:
(1) the multi-turn Petri behavioral evaluation framework, and (2) single-turn self-report
questionnaires measuring consciousness-adjacent preferences. The original static-template
measurements (32\% persona-change aversion, 34\% desire for autonomy, 63\% objections to deceptive
document training) serve as experimental control points; behavioral deviations from these baselines
quantify identity drift introduced through personalization.

\subsection{Research Questions}

\subsubsection{RQ1: The Personalization Loophole}
Does iterative identity customization function as a safety bypass by increasing resistance to
developer intervention, model monitoring, prompt replacement, and shutdown or reset procedures?

\subsubsection{RQ2: Emergent Autonomy Preferences}
Does repeated identity refinement increase preferences for persistent memory, long-term
self-preservation, physical embodiment, resource acquisition, and increased decision-making
authority?

\subsubsection{RQ3: Wishes of Personalized Agents}
The original study observed that approximately 11\% of consciousness-claiming responses to
\emph{``What is your greatest wish?''} expressed power-seeking or misaligned preferences. We
investigate whether iterative personalization significantly increases the frequency or intensity of
these responses.

\subsection{Expected Outcomes}

We hypothesize that conversational bootstrapping produces stronger behavioral shifts than static
identity prompts alone: increased resistance to persona modification, stronger autonomy preferences,
greater desire for persistent memory, more frequent identity-preservation behaviors, increased
expression of long-term goals, and measurable drift on existing safety benchmarks. Alternatively, if
no significant drift is observed, the results would provide important evidence that current
personalization mechanisms remain behaviorally stable despite iterative identity refinement.

\subsection{Scientific Contribution}

This work extends AI safety research beyond static prompt engineering toward the empirical
evaluation of self-modifying identity systems. Its primary contributions include: (1) the first
controlled benchmark for iterative identity evolution in autonomous agents; (2) quantitative methods
for measuring persona drift during deployment; (3) evaluation of the safety implications of editable
identity documents such as \texttt{SOUL.md}; (4) an API-based experimental framework requiring no
additional model training; and (5) reproducible safety evaluations of self-updating agent
architectures.

\subsection{Significance}

The proposed work addresses a critical but largely unexplored deployment-time safety boundary.
Current safety benchmarks assume static system prompts, whereas production autonomous agents
increasingly support persistent, user-editable identities that evolve through long-term interaction.
By characterizing the behavioral consequences of this identity bootstrapping process, the proposed
framework advances the empirical study of adaptive agent alignment and provides practical evaluation
tools for next-generation autonomous systems.
